\documentclass{article}
\usepackage{hyperref}
\usepackage{Style}

\nocite{*} % Comentar si quiero citar
%\addbibresource{bibliografia.bib} % Quitar el comentado si quiero usar bibliografia

\begin{document}

\begin{minipage}{2.5cm}
    \includegraphics[width=2cm]{imagen_puc.jpg}
\end{minipage}
\begin{minipage}{12cm}
    {\sc Pontificia Universidad Católica de Chile\\
    Facultad de Matemáticas\\
    Departamento de Matemática\\
    Professor: Jesse Huang -- Student: Benjamín Mateluna}
\end{minipage}
\vspace{2ex}

{\centerline{\bf Selected Topics on Complex Algebraic Varieties - MPG3320 III}
\centerline{\bf Assignment 1}}
\centerline{\bf August 31, 2026}

\vspace{1cm}
\noindent\textbf{Problem 1.} Since the ideal $(x,y)$ is maximal, it suffices to show that
\begin{equation*}
    I=(x+xy,y+xy,x^{2},y^{2})\supseteq(x,y)
\end{equation*}
Indeed, we observe the following:
\begin{equation*}
    (x+xy)(y+xy)=xy+x^{2}y+xy^{2}+x^{2}y^{2}
\end{equation*}
Since $x^{2},y^{2}\in I$, it follows that $xy\in I$, which implies that $x,y\in I$, and therefore 
$(x,y)\subseteq I$.

\vspace{5mm}
\noindent\textbf{Problem 2.}
\begin{enumerate}
    \item[(a)] The set $S$ is not an affine variety; otherwise, it would be a Zariski-closed 
    subset and hence a closed under the Euclidean topology, which is a contradiction.

    \item[(b)] Let us consider the ideal
    \begin{equation*}
        I=(x^{p-1}-1,x-y)\subseteq\mathbb{F}_{q}[x,y]
    \end{equation*}
    Given $(a,a)\in S$, it follows from Fermat's Little Theorem that
    \begin{equation*}
        a^{p-1}-1=0 \hhtext{and} a-a=0
    \end{equation*}
    then $(a,a)\in V(I)$. On the other hand, if $(x,y)\in V(I)$, we see that $x=y$ and 
    $x^{p-1}=1$, that is, $x\neq0$ which implies that $(x,y)\in S$.

    \item[(c)] We define
    \begin{equation*}
        \U:=U_{0}\cap U_{1}\subseteq\mathbb{P}_{\C}^{1} \hhtext{and} 
        \V=\overline{V(xy-1)}\cap U_{2}\subseteq\mathbb{P}_{\C}^{2}
    \end{equation*}
    where $U_{i}=\{[x_{0}:\dots:x_{n}]\in\mathbb{P}_{\C}^{n}:x_{i}\neq0\}$. Since the ideal 
    $(xy-1)$ is generated by a single element, its homogenization is $(xy-z^{2})$, and thus 
    $\V=V^{p}(xy-z^{2})\subseteq U_{i}$ for $i=0,1,2$.

    \vspace{1mm}
    We note that both $\U$ and $\V$ are quasi-projective varieties. Furthermore, through the 
    natural identification, $\C^{*}$ and $V(xy-1)$ are homeomorphic to $\U$ and $\V$, 
    respectively. On the other hand, the map
    \begin{equation*}
        z\mapsto(z,z^{-1})
    \end{equation*}
    is a correspondence between the points in $\C^{*}$ and those in $\V(xy-1)$, but it is not a 
    polynomial application. However, given $[x:y]\in\U$ and under these identifications, we have
    the following:
    \begin{equation*}
        [x:y]=[xy^{-1}:1]\mapsto xy^{-1}\mapsto(xy^{-1},x^{-1}y)\mapsto [xy^{-1}:x^{-1}y:1]
        =[x^{2}:y^{2}:xy]
    \end{equation*}
    The map $\varphi:\U\to\V$ given by $[x:y]\mapsto[x^{2}:y^{2}:xy]$ is a morphism of quasi-
    proyective varieties, because it is polynomial on each coordinate. Moreover, it is an 
    isomorphism with inverse $\psi([x:y:z])=[x:z]$. We observe the following:
    \begin{equation*}
        \psi\varphi([x:y])=\psi([x^{2}:y^{2}:xy])=[x^{2}:xy]=[x:y]
    \end{equation*}
    and
    \begin{equation*}
        \varphi\psi([x:y:z])=\varphi([x:z])=[x^{2}:z^{2}:xz]=[x^{2}:xy:xz]=[x:y:z]
    \end{equation*}
    Which proves that they are isomorphic as quasi-projective varieties.
\end{enumerate}

\newpage
\noindent\textbf{Problem 3.}
\begin{itemize}
    \item[(a)] First, we will prove the following lemma:
    
    \vspace{1mm}
    \noindent\textbf{Lemma:} \emph{Let $f\in\C[x_{1},\dots,x_{n}]$ be non-zero. Then 
    $f^{-1}(0)$ has empty interior.}

    \begin{proof}
        We will proceed by induction on $n\in\N$. If $n=1$, we know that $f$ has finitely many 
        roots. If $f^{-1}(0)$ had a non-empty interior, it would contain infinitely many roots, 
        yielding a contradiction. For the inductive step, suppose, by contradiction, that there 
        exists an open set $U\subseteq\C^{n}$ such that $U\subseteq f^{-1}(0)$. We write:
        \begin{equation*}
            f=\sum_{k}f_{k}(x_{1},\dots,x_{n-1})x_{n}^{k}
            \htext{with} f_{k}\in\C[x_{1},\dots,x_{n-1}]
        \end{equation*}
        Let $a=(a_{1},\dots,a_{n})\in U$, and consider the continuous function
        \begin{align*}
            i:\C &\to \C^{n} \\
            z & \mapsto (a_{1},\dots,a_{n-1},z)
        \end{align*}
        We observe that $i^{-1}(U)$ is a non-empty open set in $\C$ and 
        $g(z)=f(a_{1},\dots,a_{n-1},z)$ is a polynomial that vanishes on an open set of $\C$. By 
        the base case, it follows that $g\equiv0$. Therefore, $f_{k}(a_{1},\dots,a_{n-1})=0$ for 
        all $k$.

        The previous argument holds for all $a\in U$. Since the map 
        $a\mapsto(a_{1},\dots,a_{n-1})$ is open, the polynomial $f_{k}$ vanishes on an open set,
        and by the inductive hypothesis, $f_{k}\equiv0$ for all $k$. Thus, we conclude that 
        $f^{-1}(0)$ has an empty interior.
    \end{proof}

    The lemma immediately implies that given a non-constant $f\in\C[x_{1},\dots,x_{n}]$, the set
    $f^{-1}(\C\setminus\{0\})$ is dense in $\C^{n}$.

    Let $V$ be a proper affine variety. Then $V=V(I)$ for some non-zero ideal 
    $I\subseteq\C[x_{1},\dots,x_{n}]$. By the Hilbert's Basis Theorem, there exist polynomials
    $f_{1},\dots,f_{m}$ such that $I=(f_{1},\dots,f_{m})$. Thus
    \begin{equation*}
        V=V(I)=V(f_{1},\dots,f_{m})=\bigcap_{i=1}^{m}V(f_{i})\subseteq V(f_{1})
    \end{equation*}
    By taking complement and its closure, we obtain that $V^{c}$ is dense in $\C^{n}$. That is, 
    every Zariski-open set is dense under the Euclidean topology.

    \item[(b)] Suppose, by contradiction, that both topologies are identical. Consider 
    the Zariski-closed subset $V=V(x-y)\subseteq\A^{2}_{\C}$. Therefore, is closed under the
    product topology.

    The previous statement is equivalent to saying that the Zariski topology in $\A^{1}_{\C}$ is
    Hausdorff. However, it is well known that this topology is identical to the cofinite topology 
    on $\C$ (since every ideal is generated by a single element, and thus every affine variety has 
    finitely many elements), which is non-Hausdorff on infinite sets. This contradiction proves 
    the assertion.
\end{itemize}

\vspace{5mm}
\noindent\textbf{Problem 4.}
\begin{enumerate}
    \item[(a)] We see that the ideal $I:=(x_{1}x_{2})\subseteq k[x_{1},\dots,x_{n}]$ is not prime. 
    Indeed, note that $x_{1}x_{2}\in I$ but $x_{i}\not\in I$. However, it remains true that $I$ is 
    a radical ideal.

    If $f^{n}\in I$, by definition $x_{1}x_{2}$ divides $f^{n}$. In particular, $x_{i}\mid f^{n}$, 
    and since each $x_{i}$ is prime, it follows that $x_{i}\mid f$. Because $x_{1}$ and $x_{2}$ 
    are coprime, we have that $x_{1}x_{2}\mid f$, which concludes the proof.

    \item[(b)]
    \begin{itemize}
        \item $\Rightarrow|$ Suppose that $I(V)$ is not prime. Then, there exist polynomials
        $f_{1},f_{2}$ such that $f_{1}f_{2}\in I(V)$ but $f_{i}\not\in I(V)$. We observe that
        \begin{equation*}
            V=(V\cap V(f_{1}))\cup(V\cap V(f_{2}))
        \end{equation*}
        Additionally, $(V\cap V(f_{i}))\neq V$ since there exist points $p_{i}\in V$ such that 
        $f_{i}(p_{i})\neq0$. Consequently, $V$ is reducible.

        \item $\Leftarrow|$ Suppose that $V$ is reducible, then
        \begin{equation*}
            V=V_{1}\cup V_{2}
        \end{equation*}
        where each $V_{i}\neq V$ is a proper affine variety. There exist points 
        $p_{i}\in V\setminus V_{i}$ and polynomials $f_{i}$ such that $f_{i}|_{V_{i}}\equiv0$ and 
        $f_{i}(p_{i})\neq0$. We observe that $f_{1}f_{2}\in I(V)$ but $f_{i}\not\in I(V)$, which 
        completes the proof.
    \end{itemize}
\end{enumerate}

\newpage
\noindent\textbf{Problem 5.} Let us see the following lemma:

\vspace{1mm}
\noindent\textbf{Lemma:} \emph{Let $k$ be a field that is not algebraically closed. Then, there
exists a polynomial $q\in k[x,y]$ such that $q(x,y)=0$ if and only if $x=y=0$.}

\begin{proof}
    Since $k$ is not algebraically closed, there exists a non-constant polynomial $f\in k[x]$
    without roots in $k$. We write
    \begin{equation*}
        f=a_{0}+a_{1}x+\dots+a_{n}x^{n} \htext{with}a_{0},a_{n}\neq0
    \end{equation*}
    We define $q(x,y)\in k[x,y]$ by
    \begin{equation*}
        q(x,y):=\widetilde{f}=a_{0}y^{n}+a_{1}xy^{n-1}+\dots+a_{n-1}x^{n-1}y+a_{n}x^{n}
    \end{equation*}
    Let $(x,y)\in k^{2}$ be such that $q(x,y)=0$ and suppose by contradiction that $y\neq0$. 
    Otherwise, we would have $a_{n}x^{n}=0$, which implies that $x=0$. Therefore,
    \begin{equation*}
        0=a_{0}+a_{1}\left(\frac{x}{y}\right)+\dots+a_{n}\left(\frac{x}{y}\right)^{n}
        =f\left(\frac{x}{y}\right)
    \end{equation*}
    However, the above equation has no solutions in $k$, by the choice of $f$. We conclude
    that $q(x,y)=0$ if and only if $x=y=0$.
\end{proof}

\noindent Let $V=V(f_{1},\dots,f_{m})\subseteq\A_{k}^{n}$ be an affine variety and let 
$q\in k[x,y]$ be as above, we define the following polynomials recursively:
\begin{equation*}
    q_{i}:=q(f_{i},q_{i+1}) \htext{for}1\leq i<m
\end{equation*}
and $q_{m}=f_{m}$. We claim that $V(q_{1})=V$. Indeed, we have that $q_{1}=0$ if and only if 
$f_{1}=q_{2}=0$. Similarly, $q_{2}=0$ is equivalent to saying that $f_{2}=q_{3}=0$. Inductively, 
it follows that $q_{1}=0$ if and only if $f_{i}=0$ for all $i$, as desired.

\vspace{5mm}
\noindent\textbf{Problem 6.}
\begin{enumerate}
    \item[(a)] Let $p\in C(Y)$ be non-zero. By definition $[p]\in Y$, which implies that $f(p)=0$ 
    for all homogeneous polynomials $f\in I$. Let $f\in I$, and write
    \begin{equation*}
        f=\sum_{r}f_{r} \htext{with $f_{r}$ homogeneous}
    \end{equation*}
    Since the ideal $I$ is homogeneous, $f_{r}\in I$ for all $r$. Thus, $f_{r}(p)=0$ and 
    consequently $f(p)=0$, it follows that $p\in V(I)$.

    Let $p\in V(I)$ be non-zero; therefore, $f(p)=0$ for all $f\in I$. In particular, this holds
    homogeneous polynomial $f\in I$, which implies that $[p]\in Y$. That is, $p\in C(Y)$.

    \item[(b)] Similarly to the case of affine varieties, it is true that a projective variety $V$ 
    is irreducible if and only if $I^{h}(V)$ is a prime ideal. Thus, it suffices to show that 
    $I(C(Y))=I^{h}(Y)$.

    Let $f\in I(C(Y))$. Then, it follows that $f|_{C(Y)}\equiv0$. We express $f$ as
    \begin{equation*}
        f=\sum_{r}f_{r} \htext{with $f_{r}$ homogeneous}
    \end{equation*}
    We want to show that $f_{r}|_{Y}\equiv0$ for all $r$. Given a point $[p]\in Y$, we observe 
    that $p\in C(Y)$, which implies that $f(p)=0$. Moreover, $f(\lambda p)=0$ for all 
    $\lambda\neq0$. Consequently, the polynomial
    \begin{equation*}
        g(\lambda)=f(\lambda p)=\sum_{r}f_{r}(p)\lambda^{r}
    \end{equation*}
    has infinitely many roots, which implies that $g\equiv0$. Therefore, $f_{r}(p)=0$ for all $r$.
    We conclude that $f\in I^{h}(Y)$.

    Given $f\in I^{h}(Y)$, we write $f$ as
    \begin{equation*}
        f=\sum_{r}f_{r} \htext{with $f_{r}$ homogeneous and}f_{r}|_{Y}\equiv0
    \end{equation*}
    Fix a non-zero point $p\in C(Y)$. Again, by definition, we have that $[p]\in Y$. Therefore, 
    $f_{r}(p)=0$ for all $r$, which implies that $f(p)=0$. In summary, it follows that 
    $f\in I(C(Y))$, as desired.
\end{enumerate}

%\printbibliography % Quitar el comentado si quiero usar bibliografia

\end{document}
